\addtocontents{toc}{\protect\Needspace{3\baselineskip}}
\originalchapter{8.821 课程详讲：重夸克与强耦合势}
本章对应第 21 讲末及第 22 讲的 Wilson 环计算. 先从端点标签解释外部颜色源, 再明确所算环的标量耦合, 最后逐步求悬垂弦和减去重夸克自能. 这里 $\mathcal T$ 是长时间, $T$ 留给温度, $T_0=1/(2\pi\alpha')$ 是弦张力.

\originalsection{从一张分离膜制造外部颜色源}
在 $N+1$ 张 D3 中分离一张探针. 膜间开弦的非对角端点标签含 $U(N)$ 的基本表示和反基本表示, 正是原本重合规范群破缺后的重 $W$ 态. 它们不是 QCD 中新增的一种基本夸克场, 但在无限重、只保留颜色运输的极限可充当外部“夸克”.

分离距离 $\delta r$ 的态质量是 $M=\delta r/(2\pi\alpha')$. 在先取 Maldacena 解耦极限时固定 $U=\delta r/\alpha'$, 有 $M=U/(2\pi)$. 再令 $U\to\infty$, 探针到达喉部 AdS 边界, 留下一条从边界伸向内部的基本弦. 这个次序区分原渐近平直坐标 $\delta r$ 与喉部能量坐标 $U$; 不能一边要求 $r\ll R$ 取近视界、一边在未解耦的同一坐标里随意令 $r\to\infty$.

在 $N\gg1$ 且只加入少量探针时, 忽略其对原 $N$ 张膜几何的回馈, 可以在固定 AdS metric 上求该弦. 所得到的线算符由外部态如何耦合边界场决定, 不能只看一个画在平面上的路径就认定所有内部边界条件相同.

\originalsection{普通 Wilson 与 Wilson--Maldacena 环}
取连接形式 $D_\mu=\partial_\mu-\ii A_\mu$, 则普通规范 Wilson 环为 $\operatorname{Tr}\mathcal P\exp(\ii\oint A_\mu\dd x^\mu)$. 在 $\mathcal N=4$ 中, 从分离 D3 制得的重态还耦合位置标量. 常用欧氏 Wilson--Maldacena 环是
\[
W[C,n]=\frac1N\operatorname{Tr}\mathcal P
\exp\int\dd s\,[\ii A_\mu\dot x^\mu+
|\dot x|n_I\varphi^I],\qquad n_In_I=1.
\]
$n_I$ 指向分离的六维方向, 对应弦端在 $S^5$ 的位置. Lorentzian 世界线及不同标量/规范场归一化的写法可有相应 $\ii$ 和 $g_{\rm YM}$ 因子, 必须随 Wick 转动一起处理. 本章用标准欧氏固定 $S^5$ 数据的环来算经典面积; 纯规范 Wilson 的内部边界条件不同, 不将它们在全部量子修正上不加区分地视为同一个算符.

边界路径 $C$ 是弦世界面边界 $\partial\Sigma=C$, 所以在对偶处方中 $\langle W\rangle$ 由相应世界面路径积分给出. $N\to\infty$ 压低分裂/合并, $\lambda\gg1$ 使面积权重 $R^2/\alpha'=\sqrt\lambda$ 大, 领先项由经典极小面积控制. 固定 $N$ 或小 $\lambda$ 时, 这份单鞍点面积不是完整 Wilson 值.

\originalsection{一条直弦的自能}
对静态边界点, 在 Poincaré AdS 取世界面 $t,z$, 空间位置不变. 诱导欧氏 metric 为 $R^2(\dd t_E^2+\dd z^2)/z^2$, 每单位边界时间的面积能量是
\[
M(\epsilon)=\frac{R^2}{2\pi\alpha'}
\int_\epsilon^\infty\frac{\dd z}{z^2}
=\frac{\sqrt\lambda}{2\pi\epsilon}.
\]
它与重态从 $r=0$ 拉到 $r_\Lambda=R^2/\epsilon$ 的张力能 $r_\Lambda/(2\pi\alpha')$ 相同. 发散来源是外部源被指定为无限重, 不直接表示夸克对之间的相互作用无穷. 后面必须减掉两条直弦的自能, 才定义有限势.

\originalsection{矩形环为何给出静态势}
取欧氏矩形长边 $\mathcal T\gg L$, 两个外部点间距 $L$. 若 ground-state 的静态能量在长时间占优, 则裸环行为为 $\exp[-\mathcal T(2M+V(L))]$. 短边和 cusp 的局部贡献相对于长边能量的读取可另重整化. 减去 $2M\mathcal T$ 后, 正是所求相互作用势, 不把自能计入 $V$.

以 $t_E,x$ 作世界面坐标, $z=z(x)$, $z(\pm L/2)=0$, 中心反射对称. 面积为
\[
\frac{S_E}{\mathcal T}=
\frac{\sqrt\lambda}{2\pi}
\int_{-L/2}^{L/2}\dd x\,
\mathcal L(z,z'),\qquad
\mathcal L=\frac{\sqrt{1+z'^2}}{z^2}.
\]
因 $\mathcal L$ 不显含 $x$, $\mathcal L-z'\partial\mathcal L/\partial z'$ 是常数. 中心 $z=z_*,z'=0$ 固定该常数, 得
\[
\frac1{z^2\sqrt{1+z'^2}}=\frac1{z_*^2},\qquad
z'^2=\frac{z_*^4-z^4}{z^4}.
\]
弦在更深 $z_*$ 处转弯, 长距离在边界对应更深体尺度, 是径向字典的一个实际非局域例子.

\originalsection{把转向深度换成边界距离}
沿一半弦积分,
\[
\frac L2=\int_0^{z_*}\frac{z^2\dd z}{\sqrt{z_*^4-z^4}}
=z_* I_0,\quad
I_0=\int_0^1\frac{u^2\dd u}{\sqrt{1-u^4}}
=\frac{\sqrt\pi\Gamma(3/4)}{\Gamma(1/4)}.
\]
最后用 $v=u^4$ 化为 $B(3/4,1/2)/4$. 所以 $I_0\simeq0.5990701$、$z_*/L=1/(2I_0)\simeq0.8346268$. 下图使用这个实际积分生成的轮廓 CSV, 两端在 $z=0$, 顶点深度与比值一致; 图不是任意画一条看起来像悬链线的曲线.

\begin{center}
\begin{tikzpicture}
\begin{axis}[width=8.5cm,height=5cm,xlabel={$x/L$},ylabel={$z/L$},
xmin=-0.6,xmax=0.6,ymin=0,ymax=0.92,y dir=reverse,grid=major]
\addplot[structurecolor,thick] table[x=x_over_L,y=z_over_L,col sep=comma]{qa/wilson-profile.csv};
\addplot[only marks,mark=*,black] coordinates {(-0.5,0)(0.5,0)(0,0.83462684)};
\end{axis}
\end{tikzpicture}
\end{center}

\originalsection{自能扣除的有限积分}
带紫外截止的双腿能量是
\[
E_{\rm connected}=\frac{\sqrt\lambda}{\pi}
\int_\epsilon^{z_*}\frac{\dd z}
{z^2\sqrt{1-(z/z_*)^4}}.
\]
减掉 $2M=\sqrt\lambda/(\pi\epsilon)$ 后,
\[
V(L)=\frac{\sqrt\lambda}{\pi z_*}
\left[\int_0^1\frac{\dd u}{u^2}
\left(\frac1{\sqrt{1-u^4}}-1\right)-1\right].
\]
方括号的端点已经有限. 其值可由 Beta 函数有限部或分部积分求为 $-I_0$. 因而
\[
V(L)=-\frac{\sqrt\lambda I_0}{\pi z_*}
=-\frac{2I_0^2}{\pi}\frac{\sqrt\lambda}L
=-\frac{4\pi^2}{\Gamma(1/4)^4}\frac{\sqrt\lambda}L.
\]
最后用 $\Gamma(1/4)\Gamma(3/4)=\pi\sqrt2$ 检查两个常数写法一致, 数值系数约0.2284733. 势为负, 相应吸引力为
\[
F(L)=-\frac{\dd V}{\dd L}=-0.2284733\frac{\sqrt\lambda}{L^2}.
\]
负的\textbf{减自能势}并不表示裸系统能量无下界, 也不把重整化后的 Wilson 值当成概率.

共形尺度不变要求 $V\propto1/L$, 所以这项距离幂并非对偶才独有. 强耦合 $\sqrt\lambda$ 的系数则来自实际面积解; 弱耦合的领先趋势是 $\lambda/L$, 不同的耦合区间并非把 $\lambda$ 一处写错根号. 整个任意 $\lambda$ 函数仍未由经典面积求出.

\originalsection{有限温度时为何出现另一种鞍点}
把背景改为有视界的 AdS 黑膜, 连接弦无法随意越过欧氏几何的雪茄形端点. 除了一条连接面, 还可有两条各自从外部点伸到视界的面. 要比较的是真正扣除同一自能后的作用量, 不是只判断连接解是否存在.

小 $LT$ 时接近零温结果; 增大距离后, 不连通面可能先成为较小作用量, 对应静态颜色屏蔽. 连接解的存在上限与两鞍点交换优势的距离不同, 不可把它们统称一个“最大长度”. PS5.3 完整给出变分积分、数值比较、转折及所用图, 本章先说明物理判据. 下一章建立该热背景并解释热力学参数.
